\chapter{MÓDULO DE REGISTRO Y ADMINISTRACIÓN DE PARTICIPANTES}

Este módulo registra y verifica a los maestros responsables y, a través de ellos, inscribe a los estudiantes. El maestro es quien tiene cuenta: se registra con su unidad educativa, el organizador confirma con una carta firmada que de verdad enseña allí y, una vez aprobado, crea grupos para sus cursos. Los estudiantes no necesitan cuenta: reciben un código personal con el que rinden el desafío.

\section{PLANIFICACIÓN DEL MÓDULO}

Esta sección fija el objetivo del módulo, nombra sus historias de usuario y reparte el trabajo en dos iteraciones (Tabla~\ref{tab:iteraciones_participantes}).

\subsection{Objetivo del módulo}

Permitir la inscripción de unidades educativas escolares, maestros responsables y estudiantes de nivel primario y secundario, con la verificación de los maestros y grupos con código de acceso.

\subsection{Historias de usuario}

Las historias responden a tres usuarios: el maestro, que se registra e inscribe a sus estudiantes; el administrador, que revisa las solicitudes; y el estudiante, que se inscribe con el código de su grupo. Las tarjetas completas están en el Anexo A.

{\estilotabla
\begin{longtable}{l >{\raggedright\arraybackslash}p{6.9cm} >{\raggedright\arraybackslash}p{2.4cm} c c}
\arrayrulecolor{huazul}\specialrule{1.4pt}{0pt}{0pt}
\rowcolor{hufondo}\textbf{ID} & \textbf{Historia} & \textbf{Usuario} & \textbf{Prioridad} & \textbf{Esfuerzo} \\
\arrayrulecolor{hulinea}\midrule
\endhead
HU-46 & Crear una cuenta de maestro & Maestro & 1 & 8 \\
HU-47 & Confirmar el correo & Maestro & 1 & 3 \\
HU-48 & Iniciar sesión & Maestro & 1 & 3 \\
HU-49 & Conseguir la carta modelo & Maestro & 2 & 3 \\
HU-50 & Enviar los documentos de verificación & Maestro & 1 & 8 \\
HU-51 & Saber si una cuenta ya fue verificada & Maestro & 1 & 5 \\
HU-52 & Agregar otro colegio & Maestro & 2 & 5 \\
HU-53 & Ver lo que está pendiente de revisar & Administrador & 1 & 5 \\
HU-54 & Aprobar o rechazar una solicitud & Administrador & 1 & 3 \\
HU-55 & Suspender a un maestro & Administrador & 2 & 3 \\
HU-56 & Encontrar a un maestro & Administrador & 2 & 3 \\
HU-57 & Corregir los datos de un maestro & Administrador & 2 & 5 \\
HU-58 & Sumar a otro organizador & Administrador & 2 & 3 \\
HU-59 & Enlazar el sitio informativo & Administrador & 2 & 3 \\
HU-60 & Crear un grupo para un curso & Maestro & 1 & 5 \\
HU-61 & Compartir el grupo con los estudiantes & Maestro & 1 & 3 \\
HU-62 & Inscribir estudiantes uno por uno & Maestro & 1 & 8 \\
HU-63 & Inscribir a todo el curso con una planilla & Maestro & 2 & 8 \\
HU-64 & Corregir los datos de un estudiante & Maestro & 2 & 3 \\
HU-65 & Quitar a un estudiante o borrar un grupo & Maestro & 2 & 3 \\
HU-66 & Inscribirse con el código del grupo & Estudiante & 1 & 5 \\
\arrayrulecolor{hulinea}\midrule
\multicolumn{4}{r}{\textbf{Total de puntos}} & \textbf{95} \\
\arrayrulecolor{hulinea}\bottomrule
\rowcolor{white}\caption{Historias de usuario del módulo de registro y administración de participantes}\label{tab:historias_participantes}\\
\end{longtable}
\arrayrulecolor{black}
}

\subsection{Plan de entregas y alcance del módulo}

\begin{table}[H]
\estilotabla
\begin{tabularx}{\linewidth}{l >{\raggedright\arraybackslash}X >{\raggedright\arraybackslash}p{3.6cm}}
\arrayrulecolor{huazul}\specialrule{1.4pt}{0pt}{0pt}
\rowcolor{hufondo}\textbf{Iteración} & \textbf{Alcance} & \textbf{Historias} \\
\arrayrulecolor{hulinea}\midrule
Iteración 5 (3 al 14 de agosto) & Cuenta del maestro, documentos, revisión del administrador y administración de cuentas & \mbox{HU-46} a \mbox{HU-59} \\
Iteración 6 (17 al 28 de agosto) & Grupos por categoría, inscripción uno a uno o con planilla, autoinscripción con código & \mbox{HU-60} a \mbox{HU-66} \\
\arrayrulecolor{hulinea}\bottomrule
\end{tabularx}
\arrayrulecolor{black}
\caption[Plan de entregas del módulo de participantes]{Distribución de historias del módulo de registro y administración de participantes por iteración}
\label{tab:iteraciones_participantes}
\end{table}

\section{ITERACIÓN 5}

La Iteración 5 construye la cuenta del maestro y su revisión por el administrador. Se presentan su planificación, su diseño, su codificación, sus pruebas y sus resultados.

\subsection{Planificación de la iteración}

La iteración debía dejar listo el camino del maestro hasta su aprobación: registrarse, confirmar su correo, enviar sus documentos y ser revisado. Del lado del administrador, una bandeja con lo pendiente y la administración de las cuentas.

La Tabla~\ref{tab:hus_iteracion_5} lista las historias de la iteración; sus tarjetas, con los criterios de aceptación, están en el Anexo A.

\begin{table}[H]
\estilotabla
\begin{tabularx}{\linewidth}{l >{\raggedright\arraybackslash}X >{\raggedright\arraybackslash}p{2.8cm} c c}
\arrayrulecolor{huazul}\specialrule{1.4pt}{0pt}{0pt}
\rowcolor{hufondo}\textbf{ID} & \textbf{Historia} & \textbf{Usuario} & \textbf{Prioridad} & \textbf{Esfuerzo} \\
\arrayrulecolor{hulinea}\midrule
HU-46 & Crear una cuenta de maestro & Maestro & 1 & 8 \\
HU-47 & Confirmar el correo & Maestro & 1 & 3 \\
HU-48 & Iniciar sesión & Maestro & 1 & 3 \\
HU-49 & Conseguir la carta modelo & Maestro & 2 & 3 \\
HU-50 & Enviar los documentos de verificación & Maestro & 1 & 8 \\
HU-51 & Saber si una cuenta ya fue verificada & Maestro & 1 & 5 \\
HU-52 & Agregar otro colegio & Maestro & 2 & 5 \\
HU-53 & Ver lo que está pendiente de revisar & Administrador & 1 & 5 \\
HU-54 & Aprobar o rechazar una solicitud & Administrador & 1 & 3 \\
HU-55 & Suspender a un maestro & Administrador & 2 & 3 \\
HU-56 & Encontrar a un maestro & Administrador & 2 & 3 \\
HU-57 & Corregir los datos de un maestro & Administrador & 2 & 5 \\
HU-58 & Sumar a otro organizador & Administrador & 2 & 3 \\
HU-59 & Enlazar el sitio informativo & Administrador & 2 & 3 \\
\arrayrulecolor{hulinea}\midrule
\rowcolor{white}\multicolumn{4}{r}{\textbf{Total de puntos}} & \textbf{60} \\
\arrayrulecolor{hulinea}\bottomrule
\end{tabularx}
\arrayrulecolor{black}
\caption[Historias de la Iteración 5]{Historias de la Iteración 5}
\label{tab:hus_iteracion_5}
\end{table}


\subsection{Diseño}

La cuenta del maestro pasa por cuatro estados (Figura~\ref{fig:estados_maestro}). Mientras está pendiente puede entrar y completar su perfil, pero no crear grupos. Verificar el correo habilita el acceso, no la aprobación. Una cuenta aprobada o suspendida también se puede rechazar. El diagrama se lee de izquierda a derecha: el punto negro marca el registro, cada caja es un estado y cada flecha naranja, el cambio que hace el administrador, rotulado con su acción. Una cuenta rechazada todavía puede aprobarse, y una aprobada puede suspenderse y reactivarse.

\begin{figure}[H]
\centering
\begin{tikzpicture}[ancho/.style={minimum width=2.6cm}]
\node[inicio] (ini) at (-3.1,0) {};
\node[estespera, ancho] (pendiente) at (0,0) {Pendiente};
\node[estfinal, ancho] (aprobada) at (4.8,0) {Aprobada};
\node[estalerta, ancho] (suspendida) at (9.6,0) {Suspendida};
\node[estalerta, ancho] (rechazada) at (4.8,-2.8) {Rechazada};
\draw[flecha] (ini) -- node[etq, above=2pt]{registro} (pendiente);
\draw[flechaadm] (pendiente) -- node[etq, above=2pt]{aprobar} (aprobada);
\draw[flechaadm] ([yshift=5pt]aprobada.east) -- node[etq, above=2pt]{suspender} ([yshift=5pt]suspendida.west);
\draw[flechaadm] ([yshift=-5pt]suspendida.west) -- node[etq, below=2pt]{reactivar} ([yshift=-5pt]aprobada.east);
\draw[flechaadm] (pendiente.south) |- node[etq, pos=0.25, left=3pt]{rechazar} (rechazada.west);
\draw[flechaadm] ([xshift=-5pt]rechazada.north) -- node[etq, left=3pt]{aprobar} ([xshift=-5pt]aprobada.south);
\draw[flechaadm] ([xshift=5pt]aprobada.south) -- node[etq, right=3pt]{rechazar} ([xshift=5pt]rechazada.north);
\draw[flechaadm] (suspendida.south) |- node[etq, pos=0.25, right=3pt]{rechazar} (rechazada.east);
\end{tikzpicture}
\caption[Estados de la cuenta del maestro]{Estados de la cuenta del maestro}
\label{fig:estados_maestro}
\end{figure}

\begin{itemize}
    \item \textbf{Unidad educativa.} El maestro elige su colegio del catálogo nacional de 16\,099 unidades educativas. Si no está en la lista, o si enseña en casa, indica su departamento y su ciudad.
    \item \textbf{Documentos.} Con colegio, sube la carta de autorización firmada y sellada por el director, a partir de una carta modelo; si enseña en casa, su carnet por ambos lados. Los documentos son privados: solo los ven su dueño y el administrador.
    \item \textbf{Otros colegios.} Un maestro que enseña en más de un colegio pide cada uno por separado, con su propia carta y su propia aprobación.
    \item \textbf{Identidad.} La cuenta vive en Firebase Authentication, que exige verificar el correo. El servidor del Desafío comprueba el token en cada petición y solo deja a cada maestro ver y cambiar lo suyo.
\end{itemize}

\subsection{Codificación}

El registro valida cada campo mientras se escribe, y el servidor vuelve a validarlo. Las reglas de nombres, colegio escrito a mano, correo y teléfono existen en una sola versión, copiada en el frontend y el backend, y una prueba comprueba que las dos copias sean idénticas. Los documentos se revisan por su contenido real (PDF, JPG o PNG de hasta 5 MB) y se guardan en Cloudflare R2 con un nombre aleatorio. La carta modelo se genera en PDF con los datos del maestro.

El panel del administrador tiene dos vistas: «Por revisar», con cada solicitud pendiente y sus documentos, de la más antigua a la más nueva, y «Todos los maestros», con filtros por estado y búsqueda. Desde ahí se aprueba, se rechaza, se suspende y se corrigen los datos de un maestro. El administrador también nombra a otros organizadores.

\subsection{Pruebas}

\begin{table}[H]
\estilotabla
\begin{tabularx}{\linewidth}{>{\raggedright\arraybackslash}p{3.6cm} >{\raggedright\arraybackslash}X c c}
\arrayrulecolor{huazul}\specialrule{1.4pt}{0pt}{0pt}
\rowcolor{hufondo}\textbf{Tema} & \textbf{Qué se comprueba} & \textbf{Casos} & \textbf{Resultado} \\
\arrayrulecolor{hulinea}\midrule
Registro & Colegio del catálogo, a mano o educación en casa; departamento y ciudad; correo, contraseña, nombres y teléfono, en computadora y en celular & 22 & Aprobado \\
Inicio de sesión & Errores por campo y verificación del correo & 2 & Aprobado \\
Documentos & Subida desde el perfil, rechazo de archivos falsos o grandes, carta modelo en PDF y privacidad & 6 & Aprobado \\
Revisión y estados & Bandeja «Por revisar», otros colegios, aprobación, rechazo, suspensión y corrección de datos & 9 & Aprobado \\
Administración & Organizadores, ajustes del sitio y acceso restringido a lo administrativo & 7 & Aprobado \\
\arrayrulecolor{hulinea}\bottomrule
\end{tabularx}
\arrayrulecolor{black}
\caption[Pruebas de la Iteración 5]{Pruebas de extremo a extremo de la Iteración 5. Incluye dos casos agregados después de la ejecución completa del capítulo 9, que también aprueban.}
\label{tab:pruebas_i5}
\end{table}

\subsection{Resultados}

La Figura~\ref{fig:participantes_registro} muestra el comienzo del camino del maestro. En (a) está el formulario de registro, que se completa en dos pasos. En (b) está su perfil con la cuenta pendiente: el maestro ve en qué estado está y qué documentos le faltan enviar.

\parcapturas{assets/ui/desafio/registro-maestro.jpg}{Registro del maestro}{assets/ui/desafio/perfil-maestro.jpg}{Perfil con la cuenta pendiente}{Registro y perfil del maestro}{fig:participantes_registro}

En la Figura~\ref{fig:participantes_revision} se ven las dos vistas del panel del administrador. En (a), «Por revisar» lista cada solicitud con sus documentos para aprobarla o rechazarla. En (b), «Todos los maestros» permite filtrar por estado, buscar y administrar cada cuenta.

\parcapturas{assets/ui/desafio/maestros-revisar.jpg}{Solicitudes por revisar}{assets/ui/desafio/maestros-todos.jpg}{Todos los maestros}{Revisión de maestros}{fig:participantes_revision}

\section{ITERACIÓN 6}

La Iteración 6 construye los grupos y la inscripción de los estudiantes. Se presentan su planificación, su diseño, su codificación, sus pruebas y sus resultados.

\subsection{Planificación de la iteración}

Con maestros aprobados, la iteración resolvió cómo llegan los estudiantes a un desafío: el maestro crea un grupo por categoría y lo llena, él mismo o compartiendo un código para que cada estudiante se inscriba.

La Tabla~\ref{tab:hus_iteracion_6} lista las historias de la iteración; sus tarjetas, con los criterios de aceptación, están en el Anexo A.

\begin{table}[H]
\estilotabla
\begin{tabularx}{\linewidth}{l >{\raggedright\arraybackslash}X >{\raggedright\arraybackslash}p{2.8cm} c c}
\arrayrulecolor{huazul}\specialrule{1.4pt}{0pt}{0pt}
\rowcolor{hufondo}\textbf{ID} & \textbf{Historia} & \textbf{Usuario} & \textbf{Prioridad} & \textbf{Esfuerzo} \\
\arrayrulecolor{hulinea}\midrule
HU-60 & Crear un grupo para un curso & Maestro & 1 & 5 \\
HU-61 & Compartir el grupo con los estudiantes & Maestro & 1 & 3 \\
HU-62 & Inscribir estudiantes uno por uno & Maestro & 1 & 8 \\
HU-63 & Inscribir a todo el curso con una planilla & Maestro & 2 & 8 \\
HU-64 & Corregir los datos de un estudiante & Maestro & 2 & 3 \\
HU-65 & Quitar a un estudiante o borrar un grupo & Maestro & 2 & 3 \\
HU-66 & Inscribirse con el código del grupo & Estudiante & 1 & 5 \\
\arrayrulecolor{hulinea}\midrule
\rowcolor{white}\multicolumn{4}{r}{\textbf{Total de puntos}} & \textbf{35} \\
\arrayrulecolor{hulinea}\bottomrule
\end{tabularx}
\arrayrulecolor{black}
\caption[Historias de la Iteración 6]{Historias de la Iteración 6}
\label{tab:hus_iteracion_6}
\end{table}


\subsection{Diseño}

La Figura~\ref{fig:flujo_inscripcion} muestra cómo llega un estudiante a un desafío. El maestro aprobado crea un grupo para un desafío y una categoría, y el sistema le asigna un código de seis caracteres. Desde ahí hay dos caminos: el maestro inscribe a sus estudiantes uno a uno o con una planilla, o cada estudiante se inscribe con el código del grupo. En ambos casos el estudiante termina con un código personal de ocho caracteres.

\begin{figure}[H]
\centering
\begin{tikzpicture}[ancho/.style={minimum width=3.3cm, text width=3cm}]
\node[actor, ancho] (grupo) at (0,0) {El maestro aprobado crea un grupo para un desafío y una categoría};
\node[servicio, ancho] (codigo) at (4.3,0) {Código del grupo\\(6 caracteres)};
\node[actor, ancho] (maestro) at (8.6,1.5) {El maestro inscribe uno a uno o con planilla};
\node[actor, ancho] (estudiante) at (8.6,-1.5) {El estudiante se inscribe con el código};
\node[servicio, ancho] (personal) at (12.9,0) {Código personal\\(8 caracteres)};
\draw[flecha] (grupo) -- (codigo);
\draw[flecha] (codigo) |- (maestro);
\draw[flecha] (codigo) |- (estudiante);
\draw[flecha] (maestro) -| (personal);
\draw[flecha] (estudiante) -| (personal);
\end{tikzpicture}
\caption{Inscripción de estudiantes}
\label{fig:flujo_inscripcion}
\end{figure}

El grupo pertenece a un desafío y a una categoría, y sus estudiantes solo pueden elegir los cursos de esa categoría. El código personal es la única credencial del estudiante para rendir. Las reglas principales son:

\begin{itemize}
    \item Solo se inscribe mientras la inscripción del desafío está abierta.
    \item Un participante puede ser un estudiante o una pareja, si el desafío lo permite.
    \item Un nombre no puede repetirse dentro de un mismo grupo; sí en grupos distintos, porque dos estudiantes de colegios diferentes pueden llamarse igual.
    \item La planilla se importa entera o no se importa: si una fila falla, se indica cuál y no entra ninguna.
    \item El curso de un estudiante queda fijo desde que empieza a rendir; su nombre se puede seguir corrigiendo.
\end{itemize}

\subsection{Codificación}

Cada grupo tiene su pantalla, con el código grande, el enlace para compartir y dos pestañas: Estudiantes y Resultados. En Estudiantes, un formulario fijo inscribe uno tras otro, conservando el curso y la modalidad para el siguiente, y la lista permite corregir o quitar a cada uno. La planilla Excel se descarga ya preparada para el grupo, con los cursos en una lista desplegable; al subirla, el servidor la valida completa antes de escribir, y un bloqueo impide dos importaciones a la vez en el mismo desafío.

El estudiante entra en la página de inscripción, escribe el código del grupo, elige su curso con botones y confirma. Recibe su código personal y una cuenta regresiva hasta la prueba. Si su nombre ya está inscrito, se le pide el código a su maestro: el sistema nunca devuelve el código de otro estudiante.

\subsection{Pruebas}

\begin{table}[H]
\estilotabla
\begin{tabularx}{\linewidth}{>{\raggedright\arraybackslash}p{3.6cm} >{\raggedright\arraybackslash}X c c}
\arrayrulecolor{huazul}\specialrule{1.4pt}{0pt}{0pt}
\rowcolor{hufondo}\textbf{Tema} & \textbf{Qué se comprueba} & \textbf{Casos} & \textbf{Resultado} \\
\arrayrulecolor{hulinea}\midrule
Grupos & Creación con sus errores por campo, horario heredado del desafío y aislamiento entre maestros & 4 & Aprobado \\
Inscripción manual & Validación, nombres repetidos, edición y curso fijo al empezar & 5 & Aprobado \\
Planilla & Validación completa, importación atómica, hoja correcta del archivo y una importación a la vez & 6 & Aprobado \\
Autoinscripción & Errores del código, parejas, nombres por grupo, cierre de la inscripción y códigos inexistentes & 7 & Aprobado \\
Compartir & Copiar el código y el enlace del grupo & 1 & Aprobado \\
\arrayrulecolor{hulinea}\bottomrule
\end{tabularx}
\arrayrulecolor{black}
\caption[Pruebas de la Iteración 6]{Pruebas de extremo a extremo de la Iteración 6}
\label{tab:pruebas_i6}
\end{table}

\subsection{Resultados}

La Figura~\ref{fig:participantes_grupos} muestra cómo el maestro inscribe a sus estudiantes. En (a) está la lista de sus grupos. En (b), la pantalla de un grupo con sus estudiantes inscritos, cada uno con su código personal.

\parcapturas{assets/ui/desafio/grupos-lista.jpg}{Grupos del maestro}{assets/ui/desafio/grupo-estudiantes.jpg}{Estudiantes de un grupo}{Grupos e inscripción}{fig:participantes_grupos}

El estudiante también puede inscribirse solo con el código de su grupo (Figura~\ref{fig:participantes_autoinscripcion}). En (a) elige su curso con botones. En (b) recibe su código personal y la cuenta regresiva hasta la prueba.

\parcapturas{assets/ui/desafio/entrar-curso.jpg}{Elección del curso}{assets/ui/desafio/entrar-listo.jpg}{Código personal y cuenta regresiva}{Autoinscripción del estudiante}{fig:participantes_autoinscripcion}
